\chapter{Reconstruction and the closure condition}
\label{chap:closure}

% archon:covers Gluck/Closure.lean

Given a positive radius-of-curvature weight $\rho = 1/\kappa$, the source builds
a curve by integrating the unit tangent $e^{i\theta}$ against $\rho$. This
chapter records that reconstruction and the precise condition for the resulting
curve to close up.

% SOURCE: [DeTurck-Gluck], §4, p. 9, lines 163-174 (read from references/deturck-gluck-fourvertex.txt)
% SOURCE QUOTE: "There is a unique map α: [0, 2π] → R2 which begins at the
% origin and has unit tangent vector (cos θ, sin θ) and curvature κ(θ) at the
% point α(θ) . ... α(θ) = ∫0^θ (cos θ, sin θ) dθ / κ(θ) ."
\begin{definition}\leanok
[Reconstruction curve]
  \label{def:reconstruct}
  \lean{Gluck.reconstruct}
  \uses{def:radius}
  \textit{Source: DeTurck--Gluck, §4, explicit formula for $\alpha(\theta)$.}
  For a weight $\rho : \mathbb{R} \to \mathbb{R}$, the \emph{reconstruction
  curve} is
  \[
    \alpha_\rho(\theta) \;=\; \int_0^\theta e^{i\varphi}\,\rho(\varphi)\,d\varphi
      \;\in\; \mathbb{C},
  \]
  the complex form of the source's $\alpha(\theta) = \int_0^\theta (\cos\varphi,
  \sin\varphi)\,d\varphi/\kappa(\varphi)$ with $\rho = 1/\kappa$. It begins at the
  origin and has velocity $\alpha_\rho'(\theta) = e^{i\theta}\rho(\theta)$.
\end{definition}

\begin{lemma}

\leanok
[Reconstruction of a constant weight is a circle]
  \label{lem:reconstruct_const}
  \lean{Gluck.reconstruct_const}
  \uses{def:reconstruct}
  For a constant weight $\rho \equiv r$ the reconstruction curve has the closed
  form
  \[
    \alpha_{r}(\theta) \;=\; r\,i\,\bigl(1 - e^{i\theta}\bigr),
  \]
  a parametrization (by inclination angle) of the circle of radius $r$ centred at
  $r\,i$. This is the explicit curve realizing the constant curvature
  $\kappa \equiv 1/r$, used in the constant case of \cref{thm:gluck_converse}.
\end{lemma}
\begin{proof}
  \uses{def:reconstruct}
  \leanok
  The integrand is $\varphi \mapsto e^{i\varphi}r$, with antiderivative
  $F(\theta) = r\,i\,(1 - e^{i\theta})$ (indeed $F'(\theta) = r\,i\cdot(-i
  e^{i\theta}) = r\,e^{i\theta}$ since $-i\cdot i = 1$). The fundamental theorem of
  calculus gives $\alpha_r(\theta) = F(\theta) - F(0) = F(\theta)$, as
  $F(0) = 0$.
\end{proof}

\begin{lemma}\leanok
[Velocity of the reconstruction]
  \label{lem:reconstruct_deriv}
  \lean{Gluck.hasDerivAt_reconstruct}
  \uses{def:reconstruct}
  If $\rho$ is continuous then $\alpha_\rho$ is differentiable with
  $\alpha_\rho'(\theta) = e^{i\theta}\,\rho(\theta)$ for all $\theta$.
\end{lemma}
\begin{proof}

\leanok
  The integrand $\varphi \mapsto e^{i\varphi}\rho(\varphi)$ is continuous, so the
  fundamental theorem of calculus for the interval integral gives the derivative
  of $\theta \mapsto \int_0^\theta$ as the integrand evaluated at $\theta$.
\end{proof}

\begin{lemma}\leanok
[Inclination angle and induced curvature]
  \label{lem:signed_curvature_reconstruct}
  \lean{Gluck.signedCurvature_reconstruct}
  \uses{lem:reconstruct_deriv, def:signed_curvature, def:radius, def:is_curvature_function}
  Let $\kappa$ be a curvature function (\cref{def:is_curvature_function}) that is
  moreover continuously differentiable ($\kappa \in C^1$), and set $\rho = 1/\kappa$.
  Then $\rho \in C^1$, the reconstruction curve $\alpha_\rho$ is $C^2$ and regular,
  and it realizes $\kappa$ as its signed curvature: $\kappa_{\alpha_\rho}(\theta)
  = \kappa(\theta)$ for all $\theta$, with tangent line turning monotonically
  (angle of inclination exactly $\theta$).

  \emph{Regularity hypothesis is essential.} The signed-curvature formula
  (\cref{def:signed_curvature}) reads off the second derivative $\alpha_\rho''$
  (in Lean, \texttt{deriv (deriv $\alpha_\rho$)});
  without $\kappa \in C^1$ the weight $\rho = 1/\kappa$ need not be differentiable,
  so $\alpha_\rho''$ does not exist (and the Lean \texttt{deriv} returns junk $0$),
  making the conclusion false for a merely-continuous $\kappa$. The Lean signature
  therefore carries an explicit \texttt{ContDiff ℝ 1 κ} hypothesis.
\end{lemma}
\begin{proof}

\leanok
  By \cref{lem:reconstruct_deriv}, $\alpha_\rho'(\theta) = e^{i\theta}\rho(\theta)$,
  so the unit tangent is $e^{i\theta}$ (angle of inclination exactly $\theta$)
  and the speed is $\lvert\alpha_\rho'(\theta)\rvert = \rho(\theta) > 0$. Since
  $\kappa \in C^1$ and $\kappa > 0$, $\rho = 1/\kappa \in C^1$, so
  $\alpha_\rho'(\theta) = e^{i\theta}\rho(\theta)$ is itself $C^1$ and
  $\alpha_\rho''$ exists with $\alpha_\rho''(\theta) = i e^{i\theta}\rho(\theta) +
  e^{i\theta}\rho'(\theta)$. Write $\gamma' = e^{i\theta}\rho$ with
  $\Re\gamma' = \rho\cos\theta$, $\Im\gamma' = \rho\sin\theta$, and
  $\Re\gamma'' = \rho'\cos\theta - \rho\sin\theta$,
  $\Im\gamma'' = \rho'\sin\theta + \rho\cos\theta$. The signed-curvature
  numerator is
  \[
    \Re\gamma'\,\Im\gamma'' - \Im\gamma'\,\Re\gamma''
      = \rho^2(\cos^2\theta + \sin^2\theta) = \rho^2,
  \]
  the $\rho'$ contributions cancelling, while $\lvert\gamma'\rvert^3 = \rho^3$.
  Hence $\kappa_{\alpha_\rho}(\theta) = \rho^2/\rho^3 = 1/\rho(\theta) =
  \kappa(\theta)$. Equivalently, $ds/d\theta = \rho$ gives $d\theta/ds = \kappa$,
  the source's defining relation.
\end{proof}

\begin{definition}\leanok
[Closure integrals and error vector]
  \label{def:error_vector}
  \lean{Gluck.errorVector}
  \uses{def:reconstruct}
  The \emph{error vector} of the weight $\rho$ is
  \[
    E(\rho) \;=\; \alpha_\rho(2\pi)
      \;=\; \int_0^{2\pi} e^{i\theta}\,\rho(\theta)\,d\theta
      \;\in\; \mathbb{C}.
  \]
  Its real and imaginary parts are the \emph{closure integrals} $C(\rho)$ and
  $S(\rho)$ of \cref{def:closure_integral_cos,def:closure_integral_sin} below.
\end{definition}

\begin{definition}\leanok
[Cosine closure integral]
  \label{def:closure_integral_cos}
  \lean{Gluck.closureIntegralCos}
  The \emph{cosine closure integral} of the weight $\rho$ is the real number
  \[
    C(\rho) = \int_0^{2\pi}\!\cos\theta\,\rho(\theta)\,d\theta ,
  \]
  the real part of the error vector $E(\rho)$.
\end{definition}

\begin{definition}\leanok
[Sine closure integral]
  \label{def:closure_integral_sin}
  \lean{Gluck.closureIntegralSin}
  The \emph{sine closure integral} of the weight $\rho$ is the real number
  \[
    S(\rho) = \int_0^{2\pi}\!\sin\theta\,\rho(\theta)\,d\theta ,
  \]
  the imaginary part of the error vector $E(\rho)$.
\end{definition}

% SOURCE: [DeTurck-Gluck], §4, p. 9-10, lines 178-188 (read from references/deturck-gluck-fourvertex.txt)
% SOURCE QUOTE: "The error vector E = α(2π) — α(0) measures the failure of our
% curve to close up. ... we conclude that for some h \"inside\" this loop, the
% curve αh will have error vector E(h) = 0 , and hence close up to form a smooth
% simple closed curve."
\begin{lemma}\leanok
[Error vector splits into closure integrals]
  \label{lem:error_vector_eq}
  \lean{Gluck.errorVector_eq}
  \uses{def:error_vector, def:closure_integral_cos, def:closure_integral_sin}
  Suppose $\rho$ is continuous. Then $E(\rho) = C(\rho) + i\,S(\rho)$.

  The continuity hypothesis is \emph{necessary}, not cosmetic: the split into
  real and imaginary parts uses linearity of the integral, which requires the
  pieces to be integrable. For a non-integrable $\rho$ the identity is false
  (e.g. $\rho(\theta) = 1/(\theta - \pi/2)$ on $[0, 2\pi]$: the $\mathbb{C}$-valued
  Bochner integral $E(\rho)$ is $0$ since $|e^{i\theta}\rho| = |\rho|$ is
  non-integrable, while $C(\rho) = \int \cos\theta\,\rho \ne 0$). Every call site
  supplies continuity via \cref{lem:radius_props} ($\rho = 1/\kappa$ continuous).
\end{lemma}
\begin{proof}

\leanok
  Expand $e^{i\theta} = \cos\theta + i\sin\theta$ and split the integral of the
  $\mathbb{C}$-valued integrand into real and imaginary parts; linearity of the
  interval integral gives $E(\rho) = C(\rho) + i\,S(\rho)$.
\end{proof}

\begin{lemma}\leanok
[Closure criterion]
  \label{lem:closes_iff}
  \lean{Gluck.reconstruct_closes_iff}
  \uses{lem:error_vector_eq, def:reconstruct, def:is_closed_curve, lem:radius_props}
  \textit{Source: DeTurck--Gluck, §4 (error vector zero $\Leftrightarrow$ closes up).}
  Suppose $\rho = 1/\kappa$ for a curvature function $\kappa$. The reconstruction
  curve $\alpha_\rho$ extends to a closed ($2\pi$-periodic) curve if and only if
  $E(\rho) = 0$, equivalently $C(\rho) = 0$ and $S(\rho) = 0$.
\end{lemma}
\begin{proof}

\leanok
  Because $\alpha_\rho'(\theta) = e^{i\theta}\rho(\theta)$ is $2\pi$-periodic in
  $\theta$ (both $e^{i\theta}$ and $\rho$ are), the increment of $\alpha_\rho$
  over any period equals $\alpha_\rho(2\pi) - \alpha_\rho(0) = E(\rho)$. Hence the
  $2\pi$-periodic extension is well defined (single-valued) precisely when
  $E(\rho) = 0$. By \cref{lem:error_vector_eq}, $E(\rho) = 0$ iff both $C(\rho)$
  and $S(\rho)$ vanish.
\end{proof}
